% Source: https://artofproblemsolving.com/community/c3h2234478_remainder
% Stable identity: aops-topic-2234478-q01
\textbf{(a)} Write $n=19a+3$. We need $19a+3\equiv7\pmod{53}$, i.e. $19a\equiv4\pmod{53}$. Since $19\cdot14=266=5\cdot53+1$, the inverse of $19$ mod $53$ is $14$, so
\[a\equiv14\cdot4=56\equiv3\pmod{53}.\]
Then $a=53t+3$ and $n=19(53t+3)+3=1007t+60$. So the remainder mod $1007$ is $\boxed{60}$. (Check: $60=3\cdot19+3=53+7$.)

\textbf{(b)} Since $2014=2\cdot1007$, the numbers $n=1007t+60$ leave remainder $60$ mod $2014$ when $t$ is even and $1067$ when $t$ is odd. Both happen: $n=60$ and $n=1067$ both satisfy the conditions ($1067=19\cdot56+3=53\cdot20+7$). So the remainder mod $2014$ is $\boxed{60\text{ if }n\text{ is even, and }1067\text{ if }n\text{ is odd}}$; the given information does not fix it.
